\section{Conclusion}

Ultra-short PPG ở cả Awake và Drowsy chứa tổ chức động lực học theo thời gian không được noisy pseudoperiodic null đã kiểm tra tái tạo đầy đủ. Kết quả surrogate testing cho thấy tín hiệu quan sát được có cấu trúc vượt quá noisy pseudoperiodicity đơn giản, nhưng chỉ trong phạm vi null đã kiểm tra. Vì vậy, các kết quả này không được diễn giải như bằng chứng của deterministic chaos và cũng không loại trừ các mô hình ngẫu nhiên hoặc pseudoperiodic khác có cấu trúc tương quan phức tạp hơn.

Across the Awake--Drowsy transition, drowsiness liên quan đến finite-horizon forecastability thấp hơn, diagonal recurrence organization thấp hơn tại embedding nominal và estimated local trajectory divergence thấp hơn. Các thay đổi này được phản ánh bởi Mean CC giảm, Mean NRMSE tăng, DET giảm và LLE giảm trong cấu hình phân tích chính. Quan trọng hơn, ba nhóm thuộc tính này mô tả các khía cạnh khác nhau của dynamics và không nên được quy về một trục đơn giản của ``more chaos''--``less chaos'' hoặc ``more complexity''--``less complexity''. Pattern tổng hợp phù hợp hơn với một sự tái tổ chức phối hợp nhưng đa chiều của PPG dynamics trong quá trình suy giảm vigilance.

Các hướng chính nhìn chung được duy trì qua các độ dài cửa sổ $30$--$180$~s và dưới các label definitions thay thế đã khảo sát. Tuy nhiên, độ lớn hiệu ứng và mức hỗ trợ thống kê kém ổn định hơn direction, cho thấy robustness phụ thuộc vào từng metric, độ dài dữ liệu và lựa chọn tham số. Đặc biệt, DET nhạy rõ hơn với embedding delay so với Mean CC, Mean NRMSE và LLE. Sensitivity analysis đối với unknown repeated-session dependence cũng cho thấy các headline effects giữ hướng dưới hypothetical clustering đã khảo sát, nhưng không loại bỏ bất định do participant--session linkage bị thiếu. Do đó, kết quả hiện tại hỗ trợ inference ở cấp recording session hơn là hiệu ứng ở cấp participant.

Đóng góp chính của nghiên cứu là đặc trưng hóa sự tái tổ chức động lực học của ultra-short PPG qua chuyển tiếp Awake--Drowsy bằng nhiều descriptor phi tuyến bổ sung nhau, đồng thời xác định mức độ bền và giới hạn suy luận của các finding này. Nghiên cứu hiện tại chưa xác lập causal autonomic mechanisms, chưa đánh giá giá trị bổ sung đối với conventional PPG/HRV features trong một hệ thống detection, chưa chứng minh khả năng real-time hoặc early drowsiness detection và chưa có external validation. Các nghiên cứu tiếp theo cần multimodal validation với EEG, respiration hoặc các phép đo cardiovascular/autonomic đồng thời, dữ liệu có participant--session linkage được xác định, và independent replication trên các population, sensor và protocol khác để kiểm tra tính khái quát và ý nghĩa sinh lý của các dynamical patterns quan sát được.